\uuid{v31z}
\titre{Logique}
\auteur{bodin ; croizat ; sacre}
\organisation{exo7}
\licence{CC BY-NC-SA 4.0}
\datecreate{2022-03-08}
\niveau{L1}
\module{Algèbre}
\chapitre{Logique, ensemble, raisonnement}
\sousChapitre{Logique}
\difficulte{3}
\theme{Logique}
\exo7qcmid{hilisit.logique.016}

\contenu{
\begin{question}
\qclassification{100.01}
En 1761, le mathématicien suisse Lambert, ami d'Euler, démontre l'implication $\mathcal{I}$ : "$x \in \Qq \implies \tan(x) \notin \Qq$". Il remarque ensuite que $1 = \tan(\frac{\pi}{4})$. Qu'en conclut-il ?
\begin{answers}
	\bad{D'après $\mathcal{I}$, $\tan(\frac{\pi}{4}) \notin \Qq$.}
	\bad{D'après la contraposée de $\mathcal{I}$, $\tan(\frac{\pi}{4}) \notin \Qq$.}
	\bad{D'après la contraposée de $\mathcal{I}$, $\frac{\pi}{4} \in \Qq$.}
	\good{D'après la contraposée de $\mathcal{I}$, $\frac{\pi}{4} \notin \Qq$.}
\end{answers}
\begin{explanations}
On a $\tan(\frac{\pi}{4}) = 1 \in \Qq$. On utilise la contraposée de l'implication $\mathcal{I}$ : $1 = \tan(\frac{\pi}{4}) \in \Qq \implies \frac{\pi}{4} \notin \Qq$. Ceci constituait la toute première preuve de l'irrationalité de $\pi$.
\end{explanations}
\end{question}
}
